\documentclass[10pt]{article}
\usepackage[a4paper,margin=0.84in]{geometry}
\usepackage{amsmath,amssymb,amsthm}
\usepackage{booktabs}
\usepackage{microtype}
\usepackage{hyperref}
\hypersetup{colorlinks=true,linkcolor=blue,citecolor=blue,urlcolor=blue}

\newtheorem{theorem}{Theorem}
\newtheorem{lemma}{Lemma}
\newtheorem{proposition}{Proposition}
\newtheorem{remark}{Remark}
\newcommand{\norm}[1]{\left\lVert #1\right\rVert}
\newcommand{\ip}[2]{\left\langle #1,#2\right\rangle}

\title{A Signed Max-Cut Obstruction for Multi-Prime Cancellation}
\author{Liang Wang\\
School of Mathematics and Statistics\\
Huazhong University of Science and Technology (HUST)\\
Wuhan, China}
\date{28 August 2026}

\begin{document}
\maketitle

\begin{abstract}
The preceding pairwise study found exact signed cancellation directions for
individual physical prime components.  This paper asks whether a whole prime
shell admits one compatible coefficient-sign assignment.  We formulate the
question as a signed complete-graph max-cut problem: an edge is favorable
when its signed Gram cross term becomes nonpositive.  We prove that the
all-positive $K_m$ benchmark is $\lfloor m^2/4\rfloor$, that the associated
frustration index is exactly the complement of the signed maximum, and that
vertex switching preserves the objective.  An exact-rational replay of the
inherited 18-row literal grid evaluates 1,380 shell edges and 5,727 triangles.
Seventeen rows are all-positive and attain the benchmark; one early
exponent-crossover row has three negative edges and a finite sign-only gain of
three favorable edges.  The result is a structural obstruction and a route
selection device, not a magnitude-weighted energy estimate, an arithmetic
$L^2$ theorem, or a twin-prime result.
\end{abstract}

\section{Position on the twin-prime route}

This release remains in the same literal prime-shell dynamical family.  The
earlier releases separated physical prime components and measured their Gram
coherence \cite{tpc287,tpc288,tpc289}.  Nonnegative shell weighting exposed a
coherence wall \cite{tpc290}, while TPC-291 showed that a single pair can be
improved by an exact signed Schur direction \cite{tpc291}.  TPC-292 then
showed that three pair preferences obey a cycle-parity obstruction
\cite{tpc292}.  The smallest natural next question is global: does the entire
signed shell have a useful sign assignment, or is the triangle obstruction
only a local artifact?

We answer this question at the sign layer.  This distinction matters.  The
physical norm contains the magnitudes of all Gram entries, whereas the graph
objective below gives every nonzero edge unit weight.  A sign-only gain is
therefore a diagnostic and a clue for the next weighted experiment, not a
claim of physical norm reduction.

\section{Physical Gram shell}

Let $I=I_N$ be the frozen integer interval and let $\beta$ be the frozen
rational source vector.  For an odd prime $q$, height $H$, and exponent $s$,
write
\[
 K_{H,s}(h)=\frac{H^{2s}}{(H^2+h^2)^s},\qquad
 B_q(u,t)=\mathbf 1_{q\nmid u}\mathbf 1_{q\nmid t}
 \left(\mathbf 1_{u\equiv t\pmod q}-\frac1{q-1}\right).
\]
The deleted-diagonal physical component is $g_q=A_q\beta$, where
\[
 (A_q)_{u,t}=qK_{H,s}(u-t)\mathbf 1_{u\ne t}B_q(u,t).
 \label{eq:model}
\]
For a shell $S=\{q:Q<q\leq 2Q\}$, define
\[
 d_q=\norm{g_q}^2,\qquad G_{q,r}=\ip{g_q}{g_r},qquad
 \sigma_{q,r}=\operatorname{sign}(G_{q,r}).
 \label{eq:gram}
\]
All finite calculations use exact rational arithmetic, so no floating-point
rounding is used to decide an edge sign.

\section{Signed-shell formulation}

For vertex signs $a_q\in\{-1,+1\}$, call an edge favorable when
\[
 a_q a_r\sigma_{q,r}=-1.
\]
The signed-shell value and frustration index are
\[
 F(\sigma)=\max_a\#\{(q,r):a_qa_r\sigma_{q,r}=-1\},qquad
 \operatorname{fr}(\sigma)=\binom{|S|}{2}-F(\sigma).
 \label{eq:objective}
\]
Fixing one label to $+1$ removes the global reversal symmetry without losing
an optimum.

\begin{theorem}[All-positive max-cut benchmark]
If $\sigma_{q,r}=+1$ on every edge of $K_m$, then
\[
 F(\sigma)=\left\lfloor\frac{m^2}{4}\right\rfloor.
\]
\end{theorem}

\begin{proof}
Let $r$ vertices have label $+1$.  A favorable edge is exactly an edge
crossing the induced partition, so its number is $r(m-r)$.  The identity
\[
 r(m-r)=\frac{m^2}{4}-\left(r-\frac m2\right)^2
\]
gives the upper bound $\lfloor m^2/4\rfloor$, and a partition with
$r=\lfloor m/2\rfloor$ attains it.
\end{proof}

\begin{lemma}[Frustration complement and switching]
For any finite signed graph with $E$ edges, the minimum number of
unfavorable edges is $E-F(\sigma)$.  If vertex signs $t_q$ switch the graph
to $\sigma'_{q,r}=t_q\sigma_{q,r}t_r$, then
$F(\sigma')=F(\sigma)$.
\end{lemma}

\begin{proof}
For a fixed labeling every edge is either favorable or unfavorable, hence the
unfavorable count is $E$ minus the favorable count.  Taking the minimum and
maximum gives the first assertion.  For the second, map $a_q$ to
$a'_q=t_qa_q$.  Then
$a'_qa'_r\sigma'_{q,r}=a_qa_r\sigma_{q,r}$ on every edge, so the map is a
bijection preserving the objective.
\end{proof}

\begin{remark}
The graph objective is deliberately unweighted.  The actual quadratic form
is
\[
 \norm{\sum_q c_qg_q}^2=\sum_qc_q^2d_q+
 2\sum_{q<r}c_qc_rG_{q,r},
\]
so replacing $|G_{q,r}|$ by one and restricting $c_q$ to signs can change the
optimizer.  This paper makes that gap explicit.
\end{remark}

\section{Exact finite protocol}

The certificate inherits eight $s=2$ growth rows, four $s=1$ exponent
crossover rows, and six $s=2$ source-control rows.  For each row, every
prime in $Q<q\leq2Q$ is evaluated using the frozen source and physical
operator.  The resulting Gram matrix is exact; every off-diagonal sign is
then placed in a complete signed graph.  The producer enumerates all
$2^{m-1}$ labelings, records one optimum and its count, and computes the
triangle parity census as a consistency check.

The independent checker imports only the frozen TPC-268 engine and computes
the physical outputs in source-coordinate-first order.  It reconstructs the
complete row records rather than importing the TPC-293 producer.  Parent
code/result hashes and the engine hash are recorded in both checkers.

\section{Finite atlas}

\begin{table}[ht]
\centering
\small
\caption{Global exact-rational signed-shell census.}
\label{tab:global}
\begin{tabular}{@{}lr@{}}
\toprule
quantity & count\\
\midrule
rows & 18\\
shell edges & 1,380\\
maximum favorable edges (summed) & 744\\
minimum unsatisfied edges (summed) & 636\\
signed gain over all-positive benchmark & 3\\
all-positive rows & 17\\
signed-gain rows & 1\\
triangles & 5,727\\
sign-frustrated triangles & 5,718\\
\bottomrule
\end{tabular}
\end{table}

The edge signs are all positive in 17 rows.  The sole exception is
$(N,H,Q,z,s)=(256,38,27,5,1)$, whose seven-prime shell is
$\{29,31,37,41,43,47,53\}$.  Its three negative edges are
$(29,53)$, $(31,53)$, and $(41,53)$.  Consequently its signed maximum is
15, compared with the all-positive benchmark 12; its minimum unsatisfied
count is 6 instead of 9.  The aggregate gain is therefore exactly 3 edges.

\begin{table}[ht]
\centering
\scriptsize
\caption{Representative row records.  The frustration column is the exact
minimum number of unsatisfied sign edges.}
\label{tab:rows}
\begin{tabular}{@{}lrrrrrrrr@{}}
\toprule
axis & $N$ & $H$ & $Q$ & $s$ & $|S|$ & $F$ & fr & gain\\
\midrule
growth &128&24&9 &2&3 &2 &1 &0\\
growth &192&32&16&2&5 &6 &4 &0\\
growth &256&38&27&2&7 &12&9 &0\\
growth &384&50&40&2&10&25&20&0\\
growth &512&58&60&2&13&42&36&0\\
growth &512&58&90&2&17&72&64&0\\
crossover&256&38&27&1&7 &15&6 &3\\
control &384&52&70&2&15&56&49&0\\
\bottomrule
\end{tabular}
\end{table}

The 5,727 triangles contain 5,718 sign-frustrated triangles, consistent with
the TPC-292 parity census.  In particular, the late row
$(512,58,90,5,2)$ contributes all 680 possible triangles with `+++` edge
pattern.  The shell-level optimization therefore does not turn the late
near-coherent block into a uniformly favorable cancellation direction.

\section{Interpretation and claim firewall}

The result closes one structural uncertainty and leaves the arithmetic road
unchanged.  On the declared grid, the signed graph is overwhelmingly the
all-positive graph, for which the exact max-cut benchmark is already sharp.
The exceptional crossover row is useful precisely because it is small: it
shows that sign irregularity can exist, but it supplies only a unit-weight
gain and no evidence that Gram magnitudes or source coefficients exploit it.

The exact claims are the max-cut theorem, the frustration complement, and
switching invariance.  The finite atlas is certified by exact rational
reproduction and an independent physical replay.  A growing-shell signed
graph theorem, a magnitude-weighted Rayleigh estimate, a source-restricted
coefficient-image theorem, an arithmetic $L^2$ bound, fixed-power credit, and
the full Gate-B/twin-prime endpoint remain open.  The Session-named Route-A
and Route-B evaluator files are absent from this checkout; no official
evaluator pass is claimed.

\section{Reproducibility and conclusion}

From the project directory, run:
\begin{verbatim}
export PYTHONDONTWRITEBYTECODE=1
python -B code/tpc293_signed_shell_maxcut_certificate.py --check
python -B experiments/tpc293_independent_checker.py
python -B experiments/tpc293_signed_graph_stress.py
\end{verbatim}
The canonical finite result is
\texttt{results/tpc293\_certificate.json}.  Normal and optimized runs are
required to have empty standard error and byte-identical standard output.

TPC-293 promotes the route's reusable object from a signed triangle to a
signed prime-shell graph.  The next logically minimal experiment is to put
the actual Gram magnitudes back into the objective and test whether the
finite `+3` sign anomaly survives.

\bibliographystyle{plain}
\bibliography{references}
\end{document}
